% !TEX program = xelatex
% C++ 模块化构建：Modules vs Headers 与工程化迁移
\documentclass[12pt,a4paper]{ctexbook}

% ===== 页面 =====
\usepackage[margin=2.5cm]{geometry}

% ===== 字体 =====
\usepackage{fontspec}
\usepackage{iffont}
\IfFontExistsTF{Consolas}
  {\setmonofont{Consolas}}
  {\setmonofont{DejaVu Sans Mono}[Scale=MatchLowercase]}

% ===== 颜色 =====
\usepackage[dvipsnames,svgnames,x11names]{xcolor}
\definecolor{codebg}{HTML}{F5F5F5}
\definecolor{codeframe}{HTML}{E0E0E0}
\definecolor{linenumgray}{HTML}{999999}
\definecolor{cppkeyword}{HTML}{1A1AFF}
\definecolor{cppcomment}{HTML}{2E8B2E}
\definecolor{cppstring}{HTML}{B22222}
\definecolor{cppheadbg}{HTML}{2E4057}
\definecolor{csharpkeyword}{HTML}{8B008B}
\definecolor{csharptype}{HTML}{006400}
\definecolor{bashheadbg}{HTML}{4EAA25}
\definecolor{algheadbg}{HTML}{1A3C5E}

% ===== listings =====
\usepackage{listings}

% ---- C++ ----
\lstdefinestyle{cppstyle}{%
  language=C++,
  basicstyle=\small\ttfamily,numbers=left,numberstyle=\tiny\color{linenumgray},
  frame=single,rulecolor=\color{codeframe},framesep=6pt,
  breaklines=true,tabsize=4,columns=flexible,keepspaces=true,
  backgroundcolor=\color{codebg},showstringspaces=false,
  keywordstyle=\color{cppkeyword}\bfseries,
  commentstyle=\color{cppcomment}\itshape,
  stringstyle=\color{cppstring},
  morekeywords={%
    override,final,noexcept,constexpr,consteval,constinit,
    decltype,auto,concept,requires,co_await,co_yield,co_return,
    nullptr,sizeof...,static_assert,thread_local,alignas,alignof,
    namespace,using,template,typename,virtual,explicit,friend,
    mutable,volatile,register,extern,inline,
    module,import,export,
    std,string,vector,map,unordered_map,set,unordered_set,
    unique_ptr,shared_ptr,weak_ptr,optional,variant,any,
    tuple,pair,array,span,string_view,
    cin,cout,cerr,endl,
    move,forward,swap,make_unique,make_shared,
    begin,end,size,empty,push_back,emplace_back
  },
  deletekeywords={int,char,float,double,bool,void,long,short,unsigned,signed},
  emph={%
    int,char,float,double,bool,void,long,short,unsigned,signed,
    int8_t,int16_t,int32_t,int64_t,uint8_t,uint16_t,uint32_t,uint64_t,
    size_t,ptrdiff_t,intptr_t,uintptr_t,wchar_t
  },
  emphstyle=\color{csharptype},
  escapeinside={(*@}{@*)}
}

% ---- CMake ----
\lstdefinestyle{cmakestyle}{%
  basicstyle=\small\ttfamily,numbers=left,numberstyle=\tiny\color{linenumgray},
  frame=single,rulecolor=\color{codeframe},framesep=6pt,
  breaklines=true,tabsize=2,columns=flexible,keepspaces=true,
  backgroundcolor=\color{codebg},showstringspaces=false,
  keywordstyle=\color{cppkeyword}\bfseries,
  commentstyle=\color{cppcomment}\itshape,
  stringstyle=\color{cppstring},
  morekeywords={%
    cmake_minimum_required,project,add_library,add_executable,%
    target_sources,target_link_libraries,target_include_directories,%
    target_compile_features,target_compile_options,%
    find_package,include,set,option,if,else,elseif,endif,%
    foreach,endforeach,function,endfunction,macro,endmacro,%
    FILE_SET,CXX_MODULES,INTERFACE,PUBLIC,PRIVATE,
    CMAKE_CXX_STANDARD,CMAKE_CXX_STANDARD_REQUIRED,CMAKE_BUILD_TYPE,%
    install,export,configure_package_config_file
  },
  escapeinside={(*@}{@*)}
}

% ---- Shell / 命令行 ----
\lstdefinestyle{bashstyle}{%
  basicstyle=\small\ttfamily,numbers=none,
  frame=single,rulecolor=\color{codeframe},framesep=6pt,
  breaklines=true,tabsize=4,columns=flexible,keepspaces=true,
  backgroundcolor=\color{codebg},showstringspaces=false,
  keywordstyle=\color{bashheadbg}\bfseries,
  commentstyle=\color{cppcomment}\itshape,
  stringstyle=\color{cppstring},
  morekeywords={cl,clang,gcc,g++,cmake,ninja,make,msbuild,vcpkg,conan,
                pwsh,bash,python,clangd,clang-scan-deps},
  escapeinside={(*@}{@*)}
}

\lstset{style=cppstyle}

% ===== tcolorbox（统一一次加载，避免 option clash） =====
\usepackage[most]{tcolorbox}

% ===== 代码标签 =====

\newcommand{\cpplabel}{%
  \tcbox[on line,colback=cppheadbg,colframe=cppheadbg,
    boxrule=0pt,arc=1pt,left=2pt,right=2pt,top=1pt,bottom=1pt,
    colupper=white,fontupper=\scriptsize\bfseries]{C++}}

\newcommand{\cmakelabel}{%
  \tcbox[on line,colback=cppheadbg!80!black,colframe=cppheadbg!80!black,
    boxrule=0pt,arc=1pt,left=2pt,right=2pt,top=1pt,bottom=1pt,
    colupper=white,fontupper=\scriptsize\bfseries]{CMake}}

\newcommand{\shelllabel}{%
  \tcbox[on line,colback=bashheadbg,colframe=bashheadbg,
    boxrule=0pt,arc=1pt,left=2pt,right=2pt,top=1pt,bottom=1pt,
    colupper=white,fontupper=\scriptsize\bfseries]{Shell}}

% ===== tcolorbox 提示框 =====

\newtcolorbox{guideline}[1][]{
  enhanced,breakable,
  colback=blue!5!white,colframe=blue!75!black,
  fonttitle=\bfseries,title=要点,#1,
  boxed title style={colback=blue!75!black},
  attach boxed title to top left={yshift=-2mm,xshift=2mm},
  arc=2mm,boxrule=0.8mm,
}
\newtcolorbox{compare}[1][]{
  enhanced,breakable,
  colback=green!3!white,colframe=green!50!black,
  fonttitle=\bfseries,title=对比,#1,
  boxed title style={colback=green!50!black},
  attach boxed title to top left={yshift=-2mm,xshift=2mm},
  arc=2mm,boxrule=0.8mm,
}
\newtcolorbox{warning}[1][]{
  enhanced,breakable,
  colback=red!5!white,colframe=red!75!black,
  fonttitle=\bfseries,title=常见陷阱,#1,
  boxed title style={colback=red!75!black},
  attach boxed title to top left={yshift=-2mm,xshift=2mm},
  arc=2mm,boxrule=0.8mm,
}
\newtcolorbox{note}[1][]{
  enhanced,breakable,
  colback=orange!5!white,colframe=orange!80!black,
  fonttitle=\bfseries,title=注意,#1,
  boxed title style={colback=orange!80!black},
  attach boxed title to top left={yshift=-2mm,xshift=2mm},
  arc=2mm,boxrule=0.8mm,
}
\newtcolorbox{bestpractice}[1][]{
  enhanced,breakable,
  colback=purple!5!white,colframe=purple!60!black,
  fonttitle=\bfseries,title=最佳实践,#1,
  boxed title style={colback=purple!60!black},
  attach boxed title to top left={yshift=-2mm,xshift=2mm},
  arc=2mm,boxrule=0.8mm,
}
\newtcolorbox{evidence}[1][]{
  enhanced,breakable,
  colback=gray!5!white,colframe=gray!70!black,
  fonttitle=\bfseries,title=证据,#1,
  boxed title style={colback=gray!70!black},
  attach boxed title to top left={yshift=-2mm,xshift=2mm},
  arc=2mm,boxrule=0.8mm,
}

% ===== 章节标题 =====
\usepackage{titlesec}
\usepackage{fancyhdr}
\usepackage{graphicx}
\usepackage{amssymb}
\usepackage{amsmath}
\usepackage{tabularx}
\usepackage{booktabs}
\usepackage{longtable}
\usepackage{array}
\usepackage{multirow}
\usepackage{enumitem}
\usepackage{seqsplit}
\usepackage{float}

\titleformat{\chapter}[display]
  {\normalfont\huge\bfseries\color{algheadbg}}
  {\chaptertitlename\ \thechapter}{20pt}{\Huge}
\titleformat{\section}
  {\normalfont\Large\bfseries\color{blue!70!black}}
  {\thesection}{1em}{}
\titleformat{\subsection}
  {\normalfont\large\bfseries\color{blue!60!black}}
  {\thesubsection}{1em}{}

\pagestyle{fancy}
\fancyhf{}
\fancyhead[LE,RO]{\thepage}
\fancyhead[RE]{\leftmark}
\fancyhead[LO]{\rightmark}

\setlist[itemize]{leftmargin=2em,itemsep=2pt}
\setlist[enumerate]{leftmargin=2em,itemsep=2pt}

% ===== 超链接 =====
\usepackage{hyperref}
\hypersetup{
  colorlinks=true,
  linkcolor=blue!70!black,
  urlcolor=blue!60!black,
  bookmarks=true,
  pdfauthor={hsmbanlance},
  pdftitle={C++ 模块化构建参考手册},
}

% ===== 自定义命令 =====
\newcommand{\cpp}[1]{C++#1}
\newcommand{\Fn}[1]{\texttt{#1()}}
\newcommand{\Tp}[1]{\texttt{#1}}
\newcommand{\flag}[1]{\texttt{#1}}

% ====================================================================
\begin{document}

\frontmatter

\begin{titlepage}
\centering
\vspace*{4cm}
{\Huge\bfseries C++ 模块化构建参考手册\par}
\vspace{1cm}
{\LARGE Modules vs Headers \quad 优势 \quad 成熟度困局 \quad CMake 实战\par}
\vfill
{\large \today\par}
\end{titlepage}

\tableofcontents

\mainmatter

% ╔══════════════════════════════════════════════════════════════════╗
% ║  Part I — 头文件机制：本质与痛点                                ║
% ╚══════════════════════════════════════════════════════════════════╝
\part{头文件机制：本质与痛点}

\chapter{\#include 的工作模型}

\section{预处理器视角：纯文本替换}

C++ 编译流程的第一阶段是预处理。\flag{\#include} 不是语言关键字，而是预处理指令，
其语义就是把目标文件的内容逐字符粘贴到指令所在位置。经过预处理后，编译器看到的
是一个被完全展开的、没有任何 \flag{\#include} 的翻译单元（Translation Unit, TU）。

举例：一个仅有 \texttt{int main()\{\}} 的源文件，在 MSVC 下展开
\texttt{\#include <iostream>} 之后，预处理产物的行数大约在 4 万到 7 万行之间，
具体取决于标准库实现版本。换言之，每次重新编译一个翻译单元，都要重新解析这几万行
本质上和你业务无关的代码。

\begin{compare}[title=文本替换 vs 语义导入]
\begin{itemize}
  \item \textbf{\#include}：预处理器层面的复制粘贴，编译器对原始边界一无所知。
  \item \textbf{import}：编译器层面的语义导入，导入的是已经编译好的接口快照（BMI）。
\end{itemize}
\end{compare}

\section{Include Guard 与 \#pragma once 的局限}

为了避免同一个头文件在一个 TU 中被多次展开，业界长期采用两种手段：

\cpplabel
\begin{lstlisting}[style=cppstyle]
// 方式 A: 经典 include guard
#ifndef MYLIB_FOO_H
#define MYLIB_FOO_H
// ... declarations ...
#endif

// 方式 B: 编译器扩展（被所有主流编译器接受）
#pragma once
// ... declarations ...
\end{lstlisting}

但这两者只能解决“同一 TU 内的重复展开”，无法解决跨 TU 的重复解析。
即便项目里有 200 个 .cpp 都包含了同一个头文件，编译器仍要分别解析 200 次。
预编译头（PCH）是对此的局部缓解，但 PCH 本身存在粒度过粗、跨编译器不通用、
依赖维护成本高等问题。

\begin{warning}[title=include guard 的两个常见坑]
\begin{enumerate}
  \item 不同库使用相同 guard 宏名（如 \texttt{\_H\_} 或库前缀冲突）会导致
        头文件被静默跳过，链接期才报“未声明的标识符”。
  \item \flag{\#pragma once} 在符号链接、网络盘或 bind mount 下，部分编译器
        基于 inode 判重失败，仍可能重复展开。
\end{enumerate}
\end{warning}

\section{头文件机制带来的工程痛点}

\subsection{编译时间膨胀}

C++ 标准库的设计天然“头文件密集”：模板必须放在头文件中以供实例化，
\texttt{<algorithm>}、\texttt{<vector>}、\texttt{<string>} 之间存在大量传递性
include。一个中型项目里，单个 .cpp 文件展开后超过 10 万行预处理产物并不罕见。
对增量构建来说，修改一个公共头文件会触发几乎所有依赖它的 TU 重新编译。

\subsection{宏污染}

头文件中的所有 \flag{\#define} 会泄漏到包含它的所有 TU。
经典案例是 Windows 平台 \texttt{<windows.h>} 定义的 \texttt{min}、\texttt{max}、
\texttt{ERROR}、\texttt{interface} 等宏，与 \texttt{std::min}、\texttt{std::max}
冲突，迫使开发者要么定义 \texttt{NOMINMAX}，要么写
\texttt{(std::min)(a, b)} 这种带括号的诡异语法。

\subsection{顺序依赖}

某些头文件之间存在隐式顺序要求：A 头文件假设 B 头文件已被包含。
这种依赖不在文件中显式声明，编译失败时只报“A 中的某个类型未声明”，
排查困难。

\subsection{ODR 与传递性泄漏}

C++ 的“单一定义规则”（One Definition Rule, ODR）要求在同一个程序中，
某个实体只能有一个定义。头文件机制下，\texttt{inline}、模板、\texttt{constexpr}
变量等是 ODR 例外的少数载体，而绝大多数代码必须放在 .cpp 中。
但开发者经常无意中通过头文件传递依赖：A.h 因为内部用了 \texttt{std::vector}
而 include 了 \texttt{<vector>}，B.cpp 包含 A.h 后就能直接用 vector，
看似省事，实则使 B 的真实依赖被掩盖。一旦 A.h 重构去掉 vector，
B 立刻编译失败。

\begin{note}[title=头文件机制的本质局限]
头文件机制是 1970 年代 C 语言的产物，其设计假设是“源代码很小、宏很有用、
编译速度无所谓”。C++ 把这套机制扩展到模板、内联函数、constexpr、概念之后，
预处理的成本与语义复杂性已经远超原始设计意图。Modules 的诞生正是为了
从根本上替代这套机制。
\end{note}

% ╔══════════════════════════════════════════════════════════════════╗
% ║  Part II — Modules：语法、原理与构件                            ║
% ╚══════════════════════════════════════════════════════════════════╝
\part{Modules：语法、原理与构件}

\chapter{核心语法}

C++20 引入了模块（Modules）作为头文件的语义级替代品。
模块由一个或多个模块单元（module unit）组成，分为模块接口单元
（module interface unit）和模块实现单元（module implementation unit）。

\section{最小示例：定义与导入}

\cpplabel
\begin{lstlisting}[style=cppstyle]
// math.cppm  —— 模块接口单元
export module math;          // 声明这是一个名为 math 的模块接口

export int add(int a, int b) {
    return a + b;
}

int helper(int x) {          // 没有 export，对外不可见
    return x * 2;
}

// main.cpp
import math;                 // 语义级导入
import std;                  // C++23 起允许整体导入标准库

int main() {
    return add(1, 2);        // OK
    // helper(1);            // 错误：未导出
}
\end{lstlisting}

关键点：

\begin{itemize}
  \item \flag{export module math;} 必须出现在文件最前面（在 global module fragment 之后）。
  \item \flag{export} 修饰的实体构成模块的公共接口，未修饰的实体对外不可见。
  \item \flag{import math;} 不需要任何 \flag{\#include}，也不会引入宏。
  \item 模块名是字符串字面量风格的标识符，由实现自由分配命名空间，
        不会与文件系统路径绑定。
\end{itemize}

\section{Global Module Fragment}

模块接口单元在 \flag{export module} 之前可以有一段 global module fragment，
其中只能放预处理指令（如 \flag{\#include}）和声明，不能放定义：

\cpplabel
\begin{lstlisting}[style=cppstyle]
module;                       // 开启 global module fragment

#include <cstdio>             // 仅供本模块内部使用的传统头文件
#include "legacy_helper.h"

export module my.lib;         // 模块接口开始

export void do_work() {
    std::printf("works\n");   // OK，cstdio 在 GMF 中已可见
}
\end{lstlisting}

GMF 的语义是“这些声明属于全局模块（global module），不是 my.lib 的一部分”，
它们对导入方完全不可见。这是与传统头文件互操作的关键桥梁。

\section{Private Module Fragment}

C++20 还允许在模块接口单元末尾追加 private module fragment，
其中放置不希望被任何导入方看见的实现细节：

\cpplabel
\begin{lstlisting}[style=cppstyle]
export module fast.string;

export class FastString {
public:
    void append(char c);
private:
    char buf_[64];
    int  len_;
};

module :private;             // private fragment 起点

void FastString::append(char c) {
    if (len_ < 64) buf_[len_++] = c;
}
\end{lstlisting}

修改 private fragment 中的代码不会使导入方的 BMI 失效，
因此可以加快增量构建。

\section{Module Partition：分区的接口与实现}

随着单个模块的代码量增长，可以把接口拆分为多个分区
（partition），再由主接口聚合：

\cpplabel
\begin{lstlisting}[style=cppstyle]
// 主接口 math.cppm
export module math;
export import :basic;        // 重导出 basic 分区
export import :advanced;

// 分区 basic.cppm
export module math:basic;
export int add(int, int);

// 分区 advanced.cppm
export module math:advanced;
import :basic;               // 分区之间可以互相 import
export int power(int base, int exp);

// 实现单元 math_impl.cpp（不属于接口）
module math;                 // 注意：没有 export
import :basic;
int add(int a, int b) { return a + b; }
\end{lstlisting}

分区的好处：

\begin{enumerate}
  \item 接口可以拆成多个文件，便于团队协作。
  \item 修改一个分区只触发依赖该分区的其他分区重新编译，
        而不是整个模块。
  \item 分区名（如 \flag{:basic}）只在模块内部可见，
        外部用户始终只看到一个 \flag{math} 模块。
\end{enumerate}

\section{Header Units：传统头文件的模块化包装}

C++20 提供了 header unit 机制，允许把现有头文件当作模块导入：

\cpplabel
\begin{lstlisting}[style=cppstyle]
import <vector>;             // header unit
import "my_legacy.h";        // 用户头文件也可以

int main() {
    std::vector<int> v;      // 与 #include <vector> 等效
    v.push_back(1);
}
\end{lstlisting}

Header unit 的语义介于纯 \flag{\#include} 和完整 module 之间：

\begin{itemize}
  \item 头文件本身不需要修改，编译器自动把它当作匿名模块编译并缓存。
  \item 头文件中的宏依然会泄漏到导入方（这是为了向后兼容）。
  \item 同一个头文件在一个 TU 中只能 import 一次，多次 import 等同于一次。
\end{itemize}

\begin{warning}[title=Header Units 的局限]
\begin{enumerate}
  \item 编译器必须把头文件解析为合法的模块单元，遇到
        \flag{\#pragma once} + 复杂的宏拼接时可能失败。
  \item MSVC 与 Clang 在 header unit 的支持度上差异较大，
        Clang 17 之前几乎不可用。
  \item Header unit 之间互相 import 容易形成循环依赖。
\end{enumerate}
\end{warning}

\chapter{BMI：编译模型的核心}

\section{什么是 BMI}

BMI（Binary Module Interface）是模块接口被编译器解析后生成的二进制快照，
包含模块导出实体的完整语义信息：类型、模板、函数签名、概念、命名空间结构等。
它不是目标文件（.obj/.o），不能直接链接；它也不是文本头文件，不能被
预处理器读取。BMI 是编译器私有的、版本相关的中间格式。

各编译器的 BMI 文件后缀和默认路径：

\begin{itemize}
  \item \textbf{MSVC}：\texttt{.ifc}（Interface File Container），
        默认输出到构建目录。
  \item \textbf{Clang}：\texttt{.pcm}（Precompiled Module），
        默认放在模块缓存目录。
  \item \textbf{GCC}：\texttt{.gcm}（GCC Module），
        默认在 \texttt{gcm.cache/} 子目录。
\end{itemize}

\begin{warning}[title=BMI 不可移植]
BMI 是编译器内部表示，跨编译器、跨编译器版本、甚至跨编译选项都不兼容。
你不能把 MSVC 生成的 .ifc 给 Clang 用，也不能把 Clang 18 的 .pcm 给
Clang 19 用。这是当前模块化工具链最大的工程负担之一：
任何发布二进制库的尝试都必须附带特定编译器版本的 BMI。
\end{warning}

\section{编译流程对比}

\begin{compare}[title=传统 vs 模块化构建]
\textbf{传统}：每个 .cpp 独立预处理（展开所有 \#include）→ 编译 → 链接。
公共头文件被解析 $N$ 次，$N$ 是包含它的 TU 数。

\textbf{模块化}：先编译每个模块接口生成 BMI（一次）→ 编译实现单元
（直接读 BMI，跳过文本解析）→ 链接。BMI 的解析成本是 $O(1)$
而非 $O(N)$。
\end{compare}

\section{依赖扫描：构建系统的新职责}

模块的依赖关系不再能从 \flag{\#include} 行简单解析，因为：

\begin{enumerate}
  \item \flag{import math;} 中的 \texttt{math} 是模块名，不对应任何文件路径。
  \item 模块的接口可能由多个分区文件组成，需要构建系统知道分区到文件的映射。
  \item 模块依赖必须先于编译被解析，构建系统需要按拓扑顺序生成 BMI。
\end{enumerate}

P1689R5 提案规定了构建系统与编译器之间传递模块依赖信息的格式，
MSVC 通过 \flag{/scanDependencies} 输出，Clang 通过
\flag{clang-scan-deps} 工具输出。CMake 3.28+ 内部调用这些机制
完成依赖图构建。

% ╔══════════════════════════════════════════════════════════════════╗
% ║  Part III — Modules 的优势                                      ║
% ╚══════════════════════════════════════════════════════════════════╝
\part{Modules 的优势}

\chapter{为什么 Modules 是更好的抽象}

\section{编译速度：从 $O(N)$ 到 $O(1)$}

最直观的优势是编译速度。Microsoft 在 Visual Studio 2022 17.5+ 的实测中，
对一个包含 1500 个 TU 的中型项目，把
\seqsplit{\texttt{\#include\ <vector>}},
\seqsplit{\texttt{<string>}},
\seqsplit{\texttt{<algorithm>}}
等常用头文件替换为 \flag{import std;}，
全量重建时间下降约 35\%-50\%，增量构建（修改一个常用头文件）下降更为显著。

C++23 起 \flag{import std;} 成为标准库的官方模块化接口，
\texttt{std} 模块整体只需解析一次。各编译器的进度：

\begin{itemize}
  \item \textbf{MSVC 19.34+}（VS 2022 17.4）：完整支持 \flag{import std;}，
        通过 \flag{/std:c++latest} 启用。
  \item \textbf{GCC 15}：标准库模块化的初步支持。
  \item \textbf{Clang 19+}：libc++ 提供 \texttt{std.cppm}，需要手动编译为模块。
\end{itemize}

\section{真正的封装}

模块内未 \flag{export} 的实体对外完全不可见，且默认具有模块链接性
（module linkage），不会与其他 TU 中的同名实体冲突：

\cpplabel
\begin{lstlisting}[style=cppstyle]
// util.cppm
export module util;

namespace detail {           // 不需要 export，外部完全看不见
    int counter = 0;
    void log(const char* msg);
}

export void init() {
    detail::counter = 0;
}
\end{lstlisting}

对比传统头文件：在 \texttt{util.h} 中放进 \texttt{namespace detail} 的实体，
任何包含该头文件的 TU 都能看到，并可能在链接期与另一个库的同名 detail
冲突（即便用了匿名命名空间，跨 TU 也常被误用）。模块从根本上消除了
“内部细节被外部意外依赖”的可能性。

\section{没有宏污染、顺序无关}

\flag{import} 不会把任何宏引入到当前 TU。即便模块的实现内部使用了
\flag{\#define}，导入方也完全感知不到：

\cpplabel
\begin{lstlisting}[style=cppstyle]
// legacy.cppm
module;
#include <windows.h>          // 引入 min/max 等臭名昭著的宏
export module legacy;

export void do_win32_thing();

// main.cpp
import legacy;
int x = std::min(1, 2);       // 完全 OK，min 宏没有泄漏过来
\end{lstlisting}

模块之间也没有顺序依赖：\flag{import A; import B;} 与
\flag{import B; import A;} 等价，因为每个模块都是自包含的语义单元。

\section{显式接口与自描述}

模块接口文件本身就是一份精确的“这个库对外暴露什么”的清单，
不需要再去头文件目录里翻找哪些是公共的、哪些是内部的。
配合 \flag{module :private;} 还能把实现细节进一步隐藏。
对代码审查、文档生成、IDE 智能感知都更友好。

\section{ODR 友好}

模块内未导出的实体具有模块链接性，不会与其他模块/ TU 中的同名实体
冲突，从而大幅减少 ODR 违反。模板实例化的去重规则也更清晰：
同一个模板在不同模块中实例化时，编译器可以确定它们是同一个实体。

\begin{guideline}[title=Modules 的核心价值]
\begin{enumerate}
  \item 编译模型从“每个 TU 重新解析所有依赖”变为“依赖被解析一次后复用”。
  \item 封装从“约定（namespace detail）”变为“语言级保证（不 export 就不可见）”。
  \item 接口从“文本展开”变为“语义快照（BMI）”，可被工具链精确分析。
\end{enumerate}
\end{guideline}

% ╔══════════════════════════════════════════════════════════════════╗
% ║  Part IV — 成熟慢的根本原因                                     ║
% ╚══════════════════════════════════════════════════════════════════╝
\part{成熟慢的根本原因}

\chapter{Clang 支持的长期滞后}

\section{两套 Modules 的历史包袱}

Clang 团队从 2012 年开始就在做“Clang Modules”，目标是给 C 语言的
系统头文件提供模块化机制（典型用例是 Darwin SDK 的 \texttt{module.modulemap}）。
这套机制是为 C 设计的，与 C++20 Modules 是两套独立的实现，但共用大量底层代码。

C++20 Modules 标准化后，Clang 团队需要在已有 Clang Modules 的基础上
增加 C++ 专属语义（template、concept、ODR、module linkage、partition），
工程量巨大。结果是：

\begin{itemize}
  \item Clang 15（2022 年 9 月）：初步支持 \flag{-fmodules-ts}，仅能编译
        最简单的 module，分区与 header unit 几乎不可用。
  \item Clang 16-17：修复大量崩溃，但实现仍标记为 experimental。
  \item Clang 18（2024 年 3 月）：开始落地 P1815R2（明确模块依赖图、
        模块链接性等关键语义），但仍有约 30\% 的提案条目未实现。
  \item Clang 19（2024 年 9 月）：\flag{import std;} 在 libc++ 下可用，
        但需要用户自行编译 \texttt{std.cppm} 为 BMI。
  \item Clang 20（2025 年 3 月）：进一步完善，但仍被官方文档标记为
        “experimental support”。
\end{itemize}

\begin{evidence}[title=Clang 模块支持时间线证据]
\begin{itemize}
  \item Clang 官方文档
        \texttt{clang.llvm.org/docs/StandardLibrary.html}
        长期保留 “Modules support is experimental” 字样。
  \item LLVM GitHub 上以 \texttt{[Modules]} 标签的 issue 截至 2025 年仍有
        数百条 open，其中 \texttt{C++20 Modules} 子标签下的 issue 多为
        模板特化、ADL、链接错误。
  \item P1815R2 “Clarifying C++20 Modules” 是模块语义的关键修订，
        Clang 直到 18 才开始系统性落地，比 MSVC 晚约 2-3 年。
\end{itemize}
\end{evidence}

\section{典型 Bug 类别}

Clang 当前在以下场景仍频繁出问题：

\begin{enumerate}
  \item \textbf{模板特化跨模块}：在模块 A 中声明模板，在模块 B 中提供
        特化，导入方有时会得到“特化未声明”的错误。
  \item \textbf{ADL 与 using-directive}：模块边界破坏了 ADL 的传统假设，
        Clang 在某些重载解析路径上仍然出错。
  \item \textbf{Concept 与 requires 表达式}：模块导出的 concept 在
        导入方的 requires 子句中匹配失败。
  \item \textbf{链接器层面}：模板实例化的弱符号去重在不同模块间不一致，
        导致 ODR 违反或重复符号。
  \item \textbf{Global Module Fragment 边界}：GMF 中的声明对模块内可见，
        但 Clang 偶尔错误地认为它们对外可见。
\end{enumerate}

\section{标准库模块化的进度差距}

MSVC 在 VS 2022 17.4（2022 年 11 月）就提供了开箱即用的
\flag{import std;}，而 Clang 直到 19 才在 libc++ 中提供 \texttt{std.cppm}，
且需要用户手动编译为 BMI 并配置查找路径。GCC 14 之前完全不支持
\flag{import std;}，GCC 15 才提供初步实现。

这种进度差距直接导致：

\begin{itemize}
  \item 跨平台项目无法统一启用 \flag{import std;}，必须为每个编译器
        写条件编译。
  \item CI 矩阵中 Clang/GCC 的模块化构建经常失败，开发者被迫退回
        \flag{\#include}。
  \item 教程与示例多以 MSVC 为准，Clang 用户照搬时频繁踩坑。
\end{itemize}

\chapter{LSP / IDE 支持的极度滞后}

\section{clangd：模块支持仍处实验阶段}

clangd 是当前 C++ 生态最主流的 LSP 实现，与 Clang 编译器同步演进。
但 clangd 对 C++20 Modules 的支持长期落后于编译器本身，原因有三：

\begin{enumerate}
  \item \textbf{索引架构假设}：clangd 的索引基于“TU 包含的头文件路径”，
        而模块依赖不对应路径，索引器无法直接复用现有数据结构。
  \item \textbf{BMI 增量更新}：编辑一个模块接口文件需要重新生成 BMI，
        然后通知所有导入该模块的 TU 刷新。这套增量协议在 LSP 标准中
        尚未规范化。
  \item \textbf{补全 / 跳转 / 重命名}：所有这些功能都依赖语义树，
        而模块边界使得 AST 节点的所有权变得复杂（一个名字可能属于
        另一个模块的 BMI）。
\end{enumerate}

截至 2025 年，clangd 的模块支持仍需要显式启用 \flag{--enable-config}，
并在 \texttt{.clangd} 中配置类似
\seqsplit{\texttt{CompileFlags:\ \{CompilationDatabase:\ build\}}}
的选项；同时项目要正确生成
\seqsplit{\texttt{compile\_commands.json}}，
其中包含 \flag{-fmodule-file=} 等参数。即便配置正确，常见症状包括：

\begin{itemize}
  \item 模块导入后，\flag{std::} 下的符号补全列表残缺。
  \item “跳转到定义”跳到 BMI 的临时反序列化结果，而非源文件。
  \item 修改模块接口后，IDE 中所有依赖文件仍显示陈旧错误，
        需要重启 clangd。
\end{itemize}

\begin{evidence}[title=clangd 模块支持证据]
clangd 的 GitHub 仓库（\texttt{clangd/clangd}）上有大量
\texttt{[modules]} 标签的 open issue，时间跨度从 2021 年到 2025 年。
其中编号靠前的几个核心 issue（关于 BMI 缓存、import std 支持、
跨模块跳转）至今仍处于 open 或 partial fix 状态。
\end{evidence}

\section{Visual Studio IntelliSense}

Visual Studio 的 IntelliSense 在 MSVC 后端下对模块支持相对较好，
因为 MSVC 编译器与 IDE 是同一团队开发。但仍有以下限制：

\begin{itemize}
  \item 使用 Clang 后端（Visual Studio 提供 Clang 工具集选项）时，
        IntelliSense 对模块的支持几乎为零。
  \item 模块接口文件的“大纲视图”有时不显示 export 实体的层级。
  \item 修改模块后需要等待 IntelliSense 数据库刷新，期间整个项目
        显示红色波浪线，体验较差。
\end{itemize}

\section{CLion / VS Code / Neovim}

\begin{itemize}
  \item \textbf{CLion}：基于自家解析器 + clangd，模块支持随 clangd 演进，
        2024.3 版本起提供实验性开关。
  \item \textbf{VS Code}：依赖 clangd 或 Microsoft C/C++ 扩展。
        Microsoft C/C++ 扩展使用 EDG 前端，对模块的支持仅限于 MSVC 风格，
        Clang 项目几乎不可用。
  \item \textbf{Neovim}：通过 \texttt{nvim-lspconfig} 配置 clangd，
        与 VS Code 上的 clangd 体验一致。LazyVim 的
        \texttt{lang.clangd} extra 默认不开启模块支持。
\end{itemize}

\begin{warning}[title=LSP 滞后是模块化的最大体验杀手]
即便编译器支持完善，如果 IDE 在编辑过程中频繁误报、补全失败、
跳转混乱，开发者会主动放弃模块。这是当前 Clang 工具链下
模块化的真实瓶颈：不是“能不能编译”，而是“能不能愉快地写代码”。
\end{warning}

\chapter{第三方库模块化的困境}

这是即便你既是库作者又是库使用者，仍然会遇到的根本性难题。

\section{双轨发布的成本}

库作者必须同时维护两种发布形式：

\begin{enumerate}
  \item \textbf{传统头文件 + 静态/动态库}：服务于绝大多数下游用户，
        包括使用旧编译器、未启用模块的项目、跨语言绑定（如
        \flag{extern "C"} 接口）。
  \item \textbf{模块接口（.cppm 或 .ixx）+ BMI}：服务于已启用模块的下游。
\end{enumerate}

这两套发布物的代码必须严格保持语义一致。常见的做法是：

\cpplabel
\begin{lstlisting}[style=cppstyle]
// mylib.h —— 传统头文件，仍由 #include 提供
#pragma once
namespace mylib { void do_thing(); }

// mylib.cppm —— 模块接口
export module mylib;
export namespace mylib { void do_thing(); }

// mylib.cpp —— 单一实现源
#include "mylib.h"
namespace mylib { void do_thing() { /* ... */ } }
\end{lstlisting}

这种双轨结构的代价：

\begin{itemize}
  \item 接口签名变更必须同步两处，遗漏一处就会出现“头文件用户能编译，
        模块用户报错”的奇怪 bug。
  \item CI 需要为两套发布物分别构建测试。
  \item 文档（Doxygen、Sphinx）需要识别两种接口源。
\end{itemize}

\section{模板与 inline 函数的导出困境}

C++ 模板必须在每个使用点可见完整定义，传统头文件机制下这是天然的：
模板代码直接放在 .h 中。\flag{export} 一个模板到模块时，
编译器需要在 BMI 中保留完整的模板定义，以便导入方实例化。
当前编译器在以下场景仍有问题：

\begin{itemize}
  \item 模板的偏特化、变量模板、模板模板参数。
  \item 模板内部使用了 GMF 中 \flag{\#include} 进来的类型。
  \item 概念约束的模板跨模块匹配。
\end{itemize}

\texttt{inline} 函数也类似：在模块中 \flag{export inline} 的函数，
其定义需要随 BMI 一起被导入方看见。这与传统 ODR 规则的交互仍在
P1815R2 等提案中被持续澄清。

\section{ABI 兼容性}

模块的 BMI 不可跨编译器版本，这意味着：

\begin{itemize}
  \item 发布二进制库时，必须为每个目标编译器版本提供独立的 BMI。
  \item vcpkg 的 binary caching 在模块化场景下命中率大幅下降，
        因为缓存键需要包含编译器版本、模块编译选项等额外信息。
  \item Conan 2.x 的 \texttt{package\_id} 计算尚未完全覆盖模块场景。
\end{itemize}

\section{vcpkg / Conan 的支持现状}

\begin{itemize}
  \item \textbf{vcpkg}：从 2024 年起在 ports 中陆续添加模块化支持，
        但绝大多数官方 port 仍只发布传统头文件 + 库。
        vcpkg 工具链脚本 \seqsplit{\texttt{vcpkg.cmake}}
        对 \seqsplit{\texttt{FILE\_SET\ CXX\_MODULES}}
        的传递仍有边界情况。
  \item \textbf{Conan 2}：\seqsplit{\texttt{cpp\_info.builddirs}}
        与模块查找路径的集成在 2.7+ 改善，但
        \seqsplit{\texttt{conan-center-index}}
        中的 recipe 几乎都没有提供模块发布物。
\end{itemize}

\section{知名库的进度}

\begin{itemize}
  \item \textbf{fmt}：作者 Victor Zverovich 是模块化的积极推动者，
        但 fmt 自身仍未提供官方 .cppm 接口，主要原因是“下游用户的
        编译器矩阵过于碎片化”。
  \item \textbf{spdlog}：依赖 fmt，模块化进度跟随 fmt。
  \item \textbf{Boost}：Boost 1.86 起在少数库（如 Boost.URL）提供
        实验性模块支持，主体仍为头文件。
  \item \textbf{Qt}：Qt 6.7+ 提供 \texttt{qt\_add\_modules} CMake 函数，
        但本质是 Qt 自己的元对象系统模块，不是 C++20 Modules。
        Qt 团队公开表态“在工具链稳定前不会全面切换”。
  \item \textbf{Abseil}：Google 内部已大规模使用模块化构建，
        但开源版本仍只发布头文件。
  \item \textbf{标准库}：libc++ 与 MSVC STL 提供 \texttt{std} 模块，
        libstdc++ 在 GCC 15 跟进。
\end{itemize}

\section{库作者亲历：你是作者也无能为力}

即便一个库的所有代码都掌握在你手里，模块化仍然会遇到以下“非技术”难题：

\begin{enumerate}
  \item \textbf{下游编译器版本不受你控制}：你的库可以发布模块，
        但用户的 CI 用 Clang 15、Clang 16、Clang 17、GCC 13、GCC 14、
        MSVC 19.36、MSVC 19.40 各跑一遍，每一种组合都需要你提供
        对应的 BMI，或者要求用户自己编译，体验立刻劣化。
  \item \textbf{依赖的库还没模块化}：你的库 \flag{import} 了第三方库
        \texttt{dep}，但 dep 还没发布模块。你只能用 GMF
        \flag{\#include "dep.h"}，于是你的模块接口又“污染”回了头文件世界。
  \item \textbf{用户的构建系统五花八门}：CMake 3.20、3.28、Bazel、
        Meson、xmake、自研构建脚本……\flag{FILE\_SET CXX\_MODULES}
        仅在 CMake 3.28+ 可用，其他构建系统的模块支持成熟度参差不齐。
  \item \textbf{Issue 维护成本上升}：用户报告“模块导入失败”，
        你需要先排查是编译器 bug、CMake 版本、还是自己接口的问题，
        而前两者你都无法修复，只能引导用户升级工具链。
  \item \textbf{双轨发布的同步成本}：每次发布新版本，都要确保
        头文件接口与模块接口语义一致，且 CI 同时跑两套构建。
\end{enumerate}

\begin{warning}[title=库作者的现实选择]
截至 2025 年，绝大多数库作者的实际策略是：
\begin{enumerate}
  \item 继续以传统头文件为主要发布形式。
  \item 提供实验性 \texttt{.cppm} 接口（通常作为单独构建产物），
        在文档中明确标注“需要 Clang 18+ / MSVC 19.34+ / GCC 15+”。
  \item 等下游用户中“启用模块”的比例超过某个阈值（经验值约 30\%），
        再考虑把模块接口提升为一等公民。
\end{enumerate}
\end{warning}

% ╔══════════════════════════════════════════════════════════════════╗
% ║  Part V — 实战：CMake 集成与渐进迁移                            ║
% ╚══════════════════════════════════════════════════════════════════╝
\part{实战：CMake 集成与渐进迁移}

\chapter{CMake 模块化构建}

CMake 3.28（2023 年 11 月）首次提供对 C++20 Modules 的官方支持，
核心是 \flag{FILE\_SET CXX\_MODULES} 与配套的依赖扫描机制。
CMake 3.30 与 3.31 持续修复边界情况，建议生产项目至少使用 3.30。

\section{最小可用 CMakeLists.txt}

\cmakelabel
\begin{lstlisting}[style=cmakestyle]
cmake_minimum_required(VERSION 3.30)
project(mylib VERSION 1.0 LANGUAGES CXX)

set(CMAKE_CXX_STANDARD 20)
set(CMAKE_CXX_STANDARD_REQUIRED ON)
set(CMAKE_CXX_EXTENSIONS OFF)

add_library(mylib)
add_library(mylib::mylib ALIAS mylib)

# 关键：FILE_SET CXX_MODULES 声明模块接口源
target_sources(mylib
  PUBLIC
    FILE_SET HEADERS
      BASE_DIRS include
      FILES include/mylib/legacy.h
    FILE_SET CXX_MODULES
      BASE_DIRS src
      FILES
        src/mylib.cppm
        src/util.cppm
  PRIVATE
    src/mylib_impl.cpp
)

target_include_directories(mylib
  PUBLIC
    $<BUILD_INTERFACE:${CMAKE_CURRENT_SOURCE_DIR}/include>
    $<INSTALL_INTERFACE:include>
)

target_compile_features(mylib PUBLIC cxx_std_20)
\end{lstlisting}

要点说明：

\begin{itemize}
  \item \flag{FILE\_SET CXX\_MODULES} 告诉 CMake 这些是模块接口源，
        需要先生成 BMI 再编译依赖它们的 TU。
  \item \flag{PUBLIC} 修饰意味着这些模块接口会被传递给链接 mylib 的目标，
        下游 \flag{import mylib;} 时 CMake 自动配置 BMI 查找路径。
  \item \flag{BASE\_DIRS} 决定模块名的解析基准。
  \item \flag{FILE\_SET HEADERS} 与 \flag{CXX\_MODULES} 共存，
        支持双轨发布。
\end{itemize}

\section{下游消费}

\cmakelabel
\begin{lstlisting}[style=cmakestyle]
cmake_minimum_required(VERSION 3.30)
project(myapp LANGUAGES CXX)

find_package(mylib REQUIRED)

add_executable(myapp main.cpp)
target_link_libraries(myapp PRIVATE mylib::mylib)

# main.cpp 中可直接 import mylib;
# CMake 自动把 BMI 路径传给编译器
\end{lstlisting}

\section{install 与 export：发布模块化包}

\cmakelabel
\begin{lstlisting}[style=cmakestyle]
include(GNUInstallDirs)
include(CMakePackageConfigHelpers)

install(TARGETS mylib
  EXPORT mylibTargets
  FILE_SET HEADERS DESTINATION ${CMAKE_INSTALL_INCLUDEDIR}
  FILE_SET CXX_MODULES DESTINATION ${CMAKE_INSTALL_LIBDIR}/cmake/mylib
)

install(EXPORT mylibTargets
  FILE mylibTargets.cmake
  NAMESPACE mylib::
  DESTINATION ${CMAKE_INSTALL_LIBDIR}/cmake/mylib
)

configure_package_config_file(
  cmake/mylibConfig.cmake.in
  ${CMAKE_BINARY_DIR}/mylibConfig.cmake
  INSTALL_DESTINATION ${CMAKE_INSTALL_LIBDIR}/cmake/mylib
)

install(FILES ${CMAKE_BINARY_DIR}/mylibConfig.cmake
  DESTINATION ${CMAKE_INSTALL_LIBDIR}/cmake/mylib
)
\end{lstlisting}

注意 \flag{FILE\_SET CXX\_MODULES DESTINATION} 把 .cppm 源文件
（而非 BMI）安装到目标位置。下游通过 \flag{find\_package} 拿到
源文件后，由自己的编译器重新生成 BMI——这是当前规避 BMI 不可移植
问题的标准做法。

\section{Ninja 与依赖扫描}

CMake 的模块化构建强依赖支持 P1689 依赖格式的生成器，目前最完整的是
\flag{Ninja} 与 \flag{Ninja Multi-Config}。\flag{Unix Makefiles} 与
\flag{Visual Studio} 生成器的支持仍有限制。

\shelllabel
\begin{lstlisting}[style=bashstyle]
# 推荐：Ninja Multi-Config（支持多配置 + 模块化）
cmake -G "Ninja Multi-Config" -B build -S .

# MSVC 下显式启用模块（VS 2022 17.5+）
cmake -G "Ninja" -B build -S . \
  -DCMAKE_CXX_COMPILER=cl \
  -DCMAKE_CXX_FLAGS="/std:c++latest /experimental:module"

# Clang 下
cmake -G "Ninja" -B build -S . \
  -DCMAKE_CXX_COMPILER=clang++ \
  -DCMAKE_CXX_FLAGS="-std=c++20 -fmodules -fmodule-file=mylib=build/mylib.pcm"
\end{lstlisting}

CMake 内部会调用编译器的依赖扫描接口（MSVC 的
\flag{/scanDependencies}、Clang 的 \flag{clang-scan-deps}），
生成 P1689 格式的依赖文件，再据此构建 BMI 编译的拓扑顺序。

\section{常见报错与排错}

\begin{warning}[title=CMake 模块化的典型错误]
\begin{enumerate}
  \item \textbf{“Cannot find module 'mylib'”}：通常是
        \flag{FILE\_SET CXX\_MODULES} 的 \flag{BASE\_DIRS} 配置错误，
        或下游未正确链接 \flag{mylib::mylib} 别名目标。
  \item \textbf{“Module file is out of date”}：BMI 缓存与源文件不同步。
        清理构建目录或显式 \flag{cmake --build build --target clean} 后重建。
  \item \textbf{“Generator does not support C++ modules”}：当前生成器
        不支持，改用 Ninja。
  \item \textbf{“/experimental:module is deprecated”}：MSVC 19.34+
        模块支持已脱离实验，移除该 flag 即可。
\end{enumerate}
\end{warning}

\section{vcpkg / Conan 与模块化}

vcpkg 的 manifest 模式下，可以在 \texttt{vcpkg.json} 中声明依赖；
但 vcpkg port 是否提供模块化构建仍取决于上游 recipe。
若上游未提供 .cppm，CMake 侧 \flag{find\_package} 后只能继续用
\flag{\#include}。

Conan 2 的 \texttt{CMakeDeps} 生成器在 2.7+ 起会把
\flag{target\_sources} 中的 \flag{FILE\_SET CXX\_MODULES} 信息
传递到下游 \texttt{*-config.cmake}，但 \texttt{conan-center-index}
中真正提供模块的 recipe 仍是少数。

\chapter{渐进迁移策略}

\section{迁移决策树}

\begin{guideline}[title=是否启用 Modules]
按以下顺序评估：
\begin{enumerate}
  \item 编译器矩阵是否包含 MSVC 19.34+ 或 Clang 18+ 或 GCC 15+？
        若否，暂缓。
  \item 构建系统是否为 CMake 3.30+ 且生成器为 Ninja？
        若否，先升级构建系统。
  \item LSP / IDE 是否可接受当前的模块支持成熟度？
        若团队严重依赖 IDE 智能感知，建议先在内部小项目试点。
  \item 直接依赖的第三方库是否已提供模块接口？
        若否，可以用 GMF + \flag{\#include} 过渡。
\end{enumerate}
\end{guideline}

\section{三步走方案}

\begin{bestpractice}[title=推荐的渐进路径]
\begin{enumerate}
  \item \textbf{第 1 步：Header Units 过渡}。在不动现有头文件的前提下，
        把项目内部对标准库的 \flag{\#include} 改为
        \flag{import <vector>;} 等 header unit。这一步对编译器要求最低，
        能立刻享受到“标准库只解析一次”的收益。
  \item \textbf{第 2 步：内部模块、外部头文件}。把项目内部的若干
        子系统改写为模块（\flag{module myapp.db;}、
        \flag{module myapp.net;} 等），但对外发布的接口仍保留头文件形式。
        这一步只在自家代码库内闭环，不需要协调下游。
  \item \textbf{第 3 步：双轨发布}。在内部模块化稳定后，再为库添加
        \flag{.cppm} 接口与 \flag{FILE\_SET CXX\_MODULES} 安装规则，
        让下游可以选择模块化消费。
\end{enumerate}
\end{bestpractice}

\section{与 PCH 的取舍}

预编译头（PCH）与模块在“减少重复解析”这一目标上重叠。
两者的差异：

\begin{compare}[title=PCH vs Modules]
\begin{itemize}
  \item \textbf{PCH}：粒度粗（一个 .pch 包含一组头文件），
        跨编译器不通用，修改任一被包含的头文件都会使 PCH 失效。
  \item \textbf{Modules}：粒度细（每个模块独立 BMI），
        语义级缓存，修改模块内部实现不一定使下游 BMI 失效。
\end{itemize}
\end{compare}

实际项目中两者可以共存：用 PCH 包裹那些“还没模块化但被到处包含”
的稳定第三方头文件，用模块组织自家代码。

\section{展望}

C++23 引入了 \flag{import std;}，C++26 预计会进一步规范模块的
分区粒度、ABI 稳定性与跨编译器互操作。但工具链的成熟需要时间：
LLVM 团队的模块负责人 Chuanqi Xu 在 2024 年 CppCon 上指出，
Clang 模块支持要达到“生产可用”至少还需要 2-3 个主版本周期。
在此之前，绝大多数项目的现实策略是“渐进迁移、双轨发布、
保持回退路径”。

\begin{note}[title=结语]
Modules 的设计是正确的，它解决了头文件机制几十年来的根本问题。
但一项语言特性的真正落地不仅取决于标准与编译器，还取决于
工具链（构建系统、LSP、IDE、包管理器）和生态（第三方库的集体迁移）。
当前的滞后并非设计缺陷，而是工程现实——
任何一项触及编译模型底层的变革都需要十年级别的生态磨合。
\end{note}

\end{document}
